% GENERATED FILE. Edit canonical lessons, then run npm run generate:books.
% canonical-source-hash: fnv1a64:95c76c49f155cd5f
% canonical-lessons: RU-C309-pokhorony, RU-C309-khoronit, RU-C309-smert, RU-C309-vdova, RU-C309-naslednik

\chapter{Death and mourning}
\label{ch:ru-death-and-mourning}
\hlchaptermodality{309}

\begin{chapteropening}
\textbf{By the end of this chapter:} \emph{I can talk about a funeral, burying someone, death, a widow and an heir.}

Complete the last lesson of chapter 309: I can talk about a funeral, burying someone, death, a widow and an heir.
\end{chapteropening}

\section[\texorpdfstring{\ru{похороны} (pokhoróny)}{похороны (pokhoróny)}]{\ru{похороны} (pokhoróny) — a funeral}

\label{lesson:RU-C309-pokhorony}



\begin{quote}
Before the new one: say the Russian for a degree, then the Russian for a school-leaving certificate.
\end{quote}

\subsection*{You'll want to know: \ru{похороны}}
\textbf{\ru{похороны}} (\emph{pokhoróny}) — \textquotedblleft{}a funeral\textquotedblright{}.

\begin{quote}
\ru{На} \ru{похороны} \ru{пришло} \ru{много} \ru{людей}.
\end{quote}

\begin{quote}
\emph{Na pokhoróny prishló mnógo lyudéy.}
\end{quote}

\begin{quote}
A lot of people came to the funeral.
\end{quote}

Plural only; genitive \emph{\ru{похорон}}: \emph{\ru{на} \ru{похоронах}} — at the funeral.

\begin{cousinweb}
\emph{\ru{похороны}} is built on \emph{\ru{хоронить}} \textquotedblleft{}to bury\textquotedblright{}, an inherited Slavic verb whose older sense was \textquotedblleft{}to keep, to hide away\textquotedblright{}; EtymDB lists no parent.
\end{cousinweb}

\begin{grammarlens}[title={What you've built}]
One more word for telling, explaining and describing at B1.
\end{grammarlens}

\subsection*{Your turn}
\begin{itemize}
\raggedright
  \item \emph{Say it:} \emph{pokhoróny}
  \item \emph{Say it:} \emph{pokhoróny}, once more
  \item \emph{Say it:} a sentence of your own with \emph{pokhoróny}
  \item \emph{Recall:} say the Russian for a degree, then the Russian for a school-leaving certificate, then say \emph{pokhoróny} again
\end{itemize}

\begin{tcolorbox}[breakable,colback=teal!4,colframe=teal!35!black,title={Before you move on}]
What does \ru{похороны} mean? (\textquotedblleft{}A funeral\textquotedblright{}.) Say it once more.
\end{tcolorbox}
\section[\texorpdfstring{\ru{хоронить} (khoronít)}{хоронить (khoronít)}]{\ru{хоронить} (khoronít) — to bury}

\label{lesson:RU-C309-khoronit}



\begin{quote}
Before the new one: say the Russian for a school-leaving certificate, then the Russian for a funeral.
\end{quote}

\subsection*{You'll want to know: \ru{хоронить}}
\textbf{\ru{хоронить}} (\emph{khoronít}) — \textquotedblleft{}to bury\textquotedblright{}.

\begin{quote}
\ru{Его} \ru{похоронили} \ru{рядом} \ru{с} \ru{женой}.
\end{quote}

\begin{quote}
\emph{Yevó pokhoroníli ryádom s zhenóy.}
\end{quote}

\begin{quote}
He was buried next to his wife.
\end{quote}

Imperfective; perfective \emph{\ru{похоронить}} (\emph{pokhoronít}).

\begin{cousinweb}
\emph{\ru{хоронить}} is an inherited Slavic verb whose older sense was \textquotedblleft{}to keep safe, to hide away\textquotedblright{}; EtymDB lists no parent, and it is usually traced to Proto-Slavic \emph{xorniti}.
\end{cousinweb}

\begin{grammarlens}[title={What you've built}]
One more word for telling, explaining and describing at B1.
\end{grammarlens}

\subsection*{Your turn}
\begin{itemize}
\raggedright
  \item \emph{Say it:} \emph{khoronít}
  \item \emph{Say it:} \emph{khoronít}, once more
  \item \emph{Say it:} a sentence of your own with \emph{khoronít}
  \item \emph{Recall:} say the Russian for a school-leaving certificate, then the Russian for a funeral, then say \emph{khoronít} again
\end{itemize}

\begin{tcolorbox}[breakable,colback=teal!4,colframe=teal!35!black,title={Before you move on}]
What does \ru{хоронить} mean? (\textquotedblleft{}To bury\textquotedblright{}.) Say it once more.
\end{tcolorbox}
\section[\texorpdfstring{\ru{смерть} (smert)}{смерть (smert)}]{\ru{смерть} (smert) — death}

\label{lesson:RU-C309-smert}



\begin{quote}
Before the new one: say the Russian for a funeral, then the Russian for to bury.
\end{quote}

\subsection*{You'll want to know: \ru{смерть}}
\textbf{\ru{смерть}} (\emph{smert}) — \textquotedblleft{}death\textquotedblright{}.

\begin{quote}
\ru{После} \ru{смерти} \ru{отца} \ru{он} \ru{стал} \ru{хозяином} \ru{фермы}.
\end{quote}

\begin{quote}
\emph{Pósle smérti ottsá on stal khozyáinom férmy.}
\end{quote}

\begin{quote}
After his father's death he became the owner of the farm.
\end{quote}

Feminine; genitive \emph{\ru{смерти}}.

\begin{cousinweb}
\emph{\ru{смерть}} is inherited through Old East Slavic \emph{sŭmĭrtĭ} from Proto-Slavic \emph{sŭmĭrtĭ} (EtymDB), usually explained as \emph{sŭ-} and the root of \emph{\ru{умереть}} \textquotedblleft{}to die\textquotedblright{}, PIE \emph{mer-}.
\end{cousinweb}

\begin{grammarlens}[title={What you've built}]
One more word for telling, explaining and describing at B1.
\end{grammarlens}

\subsection*{Your turn}
\begin{itemize}
\raggedright
  \item \emph{Say it:} \emph{smert}
  \item \emph{Say it:} \emph{smert}, once more
  \item \emph{Say it:} a sentence of your own with \emph{smert}
  \item \emph{Recall:} say the Russian for a funeral, then the Russian for to bury, then say \emph{smert} again
\end{itemize}

\begin{tcolorbox}[breakable,colback=teal!4,colframe=teal!35!black,title={Before you move on}]
What does \ru{смерть} mean? (\textquotedblleft{}Death\textquotedblright{}.) Say it once more.
\end{tcolorbox}
\section[\texorpdfstring{\ru{вдова} (vdová)}{вдова (vdová)}]{\ru{вдова} (vdová) — a widow}

\label{lesson:RU-C309-vdova}



\begin{quote}
Before the new one: say the Russian for to bury, then the Russian for death.
\end{quote}

\subsection*{You'll want to know: \ru{вдова}}
\textbf{\ru{вдова}} (\emph{vdová}) — \textquotedblleft{}a widow\textquotedblright{}.

\begin{quote}
\ru{Вдова} \ru{жила} \ru{одна} \ru{в} \ru{большом} \ru{доме}.
\end{quote}

\begin{quote}
\emph{Vdová zhilá odná v bolshóm dóme.}
\end{quote}

\begin{quote}
The widow lived alone in a big house.
\end{quote}

A widower is \emph{\ru{вдовец}}.

\begin{cousinweb}
\emph{\ru{вдова}} is inherited from Proto-Slavic \emph{vĭdova}, from PIE \emph{h\textsubscript{1}wid\textsuperscript{h}éwh\textsubscript{2}} \textquotedblleft{}widow\textquotedblright{} (EtymDB); English \emph{widow} and Latin \emph{vidua} are cousins.
\end{cousinweb}

\begin{grammarlens}[title={What you've built}]
One more word for telling, explaining and describing at B1.
\end{grammarlens}

\subsection*{Your turn}
\begin{itemize}
\raggedright
  \item \emph{Say it:} \emph{vdová}
  \item \emph{Say it:} \emph{vdová}, once more
  \item \emph{Say it:} a sentence of your own with \emph{vdová}
  \item \emph{Recall:} say the Russian for to bury, then the Russian for death, then say \emph{vdová} again
\end{itemize}

\begin{tcolorbox}[breakable,colback=teal!4,colframe=teal!35!black,title={Before you move on}]
What does \ru{вдова} mean? (\textquotedblleft{}A widow\textquotedblright{}.) Say it once more.
\end{tcolorbox}
\section[\texorpdfstring{\ru{наследник} (naslédnik)}{наследник (naslédnik)}]{\ru{наследник} (naslédnik) — an heir}

\label{lesson:RU-C309-naslednik}



\begin{quote}
Before the new one: say the Russian for death, then the Russian for a widow.
\end{quote}

\subsection*{You'll want to know: \ru{наследник}}
\textbf{\ru{наследник}} (\emph{naslédnik}) — \textquotedblleft{}an heir\textquotedblright{}.

\begin{quote}
\ru{Сын} \ru{стал} \ru{единственным} \ru{наследником}.
\end{quote}

\begin{quote}
\emph{Syn stal yedínstvennym naslédnikom.}
\end{quote}

\begin{quote}
The son became the only heir.
\end{quote}

A woman is \emph{\ru{наследница}}.

\begin{cousinweb}
\emph{\ru{наследник}} is built on \emph{\ru{наследовать}} \textquotedblleft{}to inherit\textquotedblright{}, from \emph{\ru{на}-} and \emph{\ru{след}} \textquotedblleft{}a track\textquotedblright{} (EtymDB), with the ending \emph{-\ru{ник}} for a person.
\end{cousinweb}

\begin{grammarlens}[title={What you've built}]
One more word for telling, explaining and describing at B1.
\end{grammarlens}

\subsection*{Your turn}
\begin{itemize}
\raggedright
  \item \emph{Say it:} \emph{naslédnik}
  \item \emph{Say it:} \emph{naslédnik}, once more
  \item \emph{Say it:} a sentence of your own with \emph{naslédnik}
  \item \emph{Recall:} say the Russian for death, then the Russian for a widow, then say \emph{naslédnik} again
\end{itemize}

\begin{tcolorbox}[breakable,colback=teal!4,colframe=teal!35!black,title={Before you move on}]
What does \ru{наследник} mean? (\textquotedblleft{}An heir\textquotedblright{}.) Say it once more.
\end{tcolorbox}
